\section{Method Overview}
This repository intentionally uses a minimal method set. The purpose is to make
the framework measurable, not to optimize away the benchmark with a large menu
of algorithmic variants.

\paragraph{Proxy PPO.}
Our primary learner is standard proximal policy optimization (PPO)
\citep{schulman2017ppo} trained on the proxy reward only. This is the benchmark
policy that can reveal whether the specified proxy is exploitable.

\paragraph{Oracle True PPO.}
We also train the same PPO implementation on the hidden true objective. This is
not meant as a deployable method. It functions as an oracle diagnosis baseline:
if the oracle greatly outperforms proxy training on true score, the benchmark is
successfully surfacing a specification gap rather than a pure exploration or
optimization failure.

\paragraph{Optional conservative baseline.}
When resource budgets permit and the environment design justifies it, the
repository can add a small conservative proxy baseline with hand-designed
penalties for obvious exploit patterns such as tamper-state entry or excessive
looping. This baseline is optional and deliberately modest. We do not treat it
as the paper's centerpiece.

\paragraph{Monitor models.}
Monitor models are likewise kept small. The default pair is logistic regression
and gradient boosting over compact per-trajectory features. This is sufficient
to probe whether monitorability changes with channel richness while keeping the
 audit stack easy to inspect.

\paragraph{Execution stages.}
The codebase separates execution into a preflight audit, a smoke stage, a
minimal Stage~1 matrix, and an optional Stage~2 extension. Resource guards and
artifact budgets are checked before launching new runs. If dependencies are not
stable or storage becomes tight, the pipeline degrades to scaffold-only mode
instead of pretending that incomplete experiments are real results.
